\section{Conclusion}
\label{sec:conclusion}

This paper has presented the \emph{Kiwitsche 2-Punkt-Verfahren} (KZV), a
practical procedure for estimating optimal PID control parameters using
Bang-Bang control as a starting point.

The method exploits the stationary limit-cycle oscillation that a Bang-Bang
controller naturally produces on a stable plant.
Using describing function analysis, three FOPDT model parameters (gain,
time constant, dead time) are identified from the measured oscillation period
$\Tu$, amplitude $\Au$, and the known relay switching levels.
These parameters are then fed into IMC-based PID tuning formulas to obtain
$\Kp$, $\Ti$, and $\Td$.

Simulation results on the benchmark FOPDT plant (\cref{sec:results})
demonstrate that KZV-IMC achieves \SI{5.5}{\percent} overshoot and an IAE
within \SI{2}{\percent} of the theoretically optimal IMC controller, while
Ziegler--Nichols tuning yields \SI{57.9}{\percent} overshoot and more than
eight times the IAE.

\Cref{sec:nonlinear} extends the analysis to a nonlinear FOPDT-like plant.
The KZV is shown to accurately recover the local linearisation at each
operating point, enabling a simple gain-scheduling strategy: one additional
Bang-Bang experiment at any new setpoint is sufficient to retune the PID
optimally without requiring an explicit nonlinear plant model.

The main advantages of the KZV over classical methods are:
\begin{itemize}
    \item \textbf{No additional experiment needed}: the Bang-Bang phase
          already present during commissioning is sufficient.
    \item \textbf{Dead-time estimation}: unlike basic ZN relay feedback,
          the KZV provides a separate estimate of dead time, enabling
          more robust IMC-style tuning.
    \item \textbf{Simple implementation}: all computations involve closed-form
          expressions and can be executed on standard PLC hardware.
\end{itemize}

\subsection*{Future Work}

Open directions include:
\begin{itemize}
    \item Experimental validation on a physical cooling test bench.
    \item Extension to second-order identification using multi-frequency
          relay tests.
    \item Online (recursive) adaptation of the PID parameters as the plant
          ages or operating conditions change.
    \item Integration with auto-tuning firmware for industrial controllers.
\end{itemize}
